% =====================================================================
% A2 — Section 5: Ablations.  DRAFT 1, 6 Oct 2026.
% Table 2 carries the numbers so the prose can stay short. Comparisons are
% seed-level sign tests (runs sharing a seed are dependent across sizes:
% same targets, initialisation and nested training sets), lambda_3 > 0 only.
% =====================================================================
\section{Ablations}
\label{sec:ablations}

Since $\lambda_3 > 0$ is not enough, something in training must decide whether
a run finds the structure. The paper
describes a ``competition between representation learning and decoder
overfitting'' (lower half of Fig.~\ref{fig:pipeline}). Decoder weight decay and the decoder
learning rate are two of the settings the paper's phase diagrams use to
control this competition, so I varied them. Table~\ref{tab:ablations}
summarises the results, and Fig.~\ref{fig:ablations} shows them per training
size.

\begin{table}[t]
  \centering
  \caption{All configurations, 220 runs each on the same training sets.
  Successes, where runs reaching $\mathrm{RQI}>0.95$ within $10^4$ steps, split by
  whether the run's training set has $\lambda_3>0$. $r$ at 50\% are training
  data fraction at which half of all runs succeed. $+$/$-$ are number of the 20 seeds whose successes with
  $\lambda_3>0$ go up/down relative to the released configuration (decoder
  learning rate $10^{-4}$, no weight decay). $p$ are exact sign test over seeds.}
  \label{tab:ablations}
  \small
  \setlength{\tabcolsep}{4pt}
  \begin{tabular}{@{}lccccc@{}}
    \toprule
    & \multicolumn{2}{c}{Successes} & & \multicolumn{2}{c}{Seeds vs.\ released} \\
    \cmidrule(lr){2-3} \cmidrule(l){5-6}
    Decoder setting & $\lambda_3{>}0$ & $\lambda_3{=}0$ & \shortstack{$r$ at\\50\%} & $+$/$-$ & $p$ \\
    \midrule
    Released & 70/138 & 0/82 & 0.76 & - & - \\
    Weight decay 1   & 71/138 & 0/82 & 0.76 & 1/0  & 1.0 \\
    Weight decay 5   & 74/138 & 0/82 & 0.76 & 4/0  & 0.13 \\
    Learning rate $10^{-5}$ & 96/138 & 0/82 & 0.62 & 14/2 & 0.004 \\
    Learning rate $10^{-3}$ & 15/138 & 0/82 & -   & 0/20 & $2{\times}10^{-6}$ \\
    \bottomrule
  \end{tabular}
\end{table}


\subsection{A1: decoder weight decay}
\label{sec:abl-wd}

Weight decay does not matter much. With decay 1 or 5, 71 and 74 of the 138 runs
with $\lambda_3 > 0$ succeed, compared with 70 without it. The point where
half the runs succeed stays at $r \approx 0.76$. With a decoder learning rate of $10^{-4}$, a weight decay of
5 in AdamW would on its own shrink the decoder weights to under 1\% of their
size over $10^4$ steps. Even so, only four of the 68 failed runs recover.

\subsection{A2: decoder learning rate}
\label{sec:abl-lr}

The decoder learning rate has a large effect, in the direction the paper's
phase diagrams suggest. Setting the decoder learning rate to $10^{-5}$ raises the
number of successful runs with $\lambda_3 > 0$ from 70 to 96, and moves the
50\% point from $r \approx 0.76$ to $r \approx 0.62$. On the same training
sets, 28 runs go from failure to success and only two the other way, and
14 of the 20 seeds improve. A
faster decoder ($10^{-3}$) does the opposite. Only 15 runs succeed, and the
success probability never gets above 0.30. Its 123 failed runs with $\lambda_3 > 0$ fit the training set (at least
93\%) but get only 25\% of the held-out pairs right.

With the slow decoder, successful runs need fewer steps as $r$ grows, where the median is about 3500 at $r = 0.64$ and 2500 at $r = 0.98$, as the paper reports. Within
one training size, steps and $\lambda_3$ show no consistent relation. The slow decoder does not close the gap completely, since 42 runs with
$\lambda_3 > 0$ still fail. I retrained eight of them for $5 \times 10^4$
steps. One passed the threshold at step 21,217 and
then fell back to 0.67, and the other seven stayed below RQI 0.4.
